\bookchapter{פרק 30: הקרנה, מצגות וחיבור לרכב}

\chapterintro{מסך הטלפון 7 אינץ', אך הטלוויזיה 55 והמקרן 100. העברת התמונה למסך גדול היא חוויה שונה. פרק זה מסביר הקרנה, מצגות וחיבור לרכב.}

\booksection{א', שלוש מערכות הקרנה אלחוטית - אל תתבלבל}

הקרנה אלחוטית נראית אותו דבר, אך יש שלוש טכנולוגיות שונות:

\begin{center}
\begin{tabularx}{\linewidth}{|X|X|X|}
\hline
\textbf{טכנולוגיה} & \textbf{מי משתמש} & \textbf{דרישות אביזרים} \tabularnewline
\hline
AirPlay & iPhone/iPad/Mac & אפל, הכי חלק. דרוש Apple TV או טלוויזיה עם AirPlay \tabularnewline
Miracast & Android & Wi-Fi Direct, ללא נתב. ברירת מחדל ב-Android \tabularnewline
DLNA & כל הפלטפורמות & דוחף מדיה (וידאו/מוזיקה/תמונות), לא משקף מסך. להקרנת וידאו \tabularnewline
\hline
\end{tabularx}
\end{center}

\booksubsection{שיטות לפי פלטפורמה:}

\begin{itemize}
  \item \textbf{אנדרואיד}:מרכז בקרה > שיקוף מסך > בחר טלוויזיה/Apple TV. אם לא תומך, קנה Apple TV או מקל.
  \item \textbf{HarmonyOS}:מסך התראות > "הקרנה אלחוטית"/"הקרנה חכמה" > בחר טלוויזיה. דרוש אותו Wi-Fi.
  \item \textbf{פתרון אוניברסלי}:אם הטלוויזיה ישנה, קנה Chromecast או ממיר, חבר ל-HDMI.
\end{itemize}

\booksection{ב', הקרנה בכבל - יציב יותר, בלי השהיה}

אלחוטי נוח אך עם השהיה, למשחק ומצגות פחות טוב. כבל פותר:

\begin{itemize}
  \item \textbf{אנדרואיד}:כבל Lightning/USB-C ל-HDMI (100-200). חבר והתמונה יוצאת ל-HDMI.
  \item \textbf{HarmonyOS}:USB-C ל-HDMI ישיר (דרוש DP Alt Mode, רוב הדגלים). כבל USB-C to HDMI.
  \item \textbf{חדר ישיבות}:אם מציג, החזק כבל HDMI בתיק. כל חדר ישיבות יש HDMI.
\end{itemize}

\booksection{ג', טלפון למחשב - מצב שולחן עבודה}

חלק ממכשירי Android תומכים ב“מצב שולחן עבודה” - חיבור לצג וממשק דמוי מחשב:

\begin{itemize}
  \item \textbf{Samsung DeX}:תכונת הדגל של Samsung. חבר צג וקבל שולחן עבודה, חלונות, קרוב למחשב. עם Bluetooth, בלי מחשב נייד.
  \item \textbf{מצב שולחן Huawei}:Huawei Mate/Pura עם צג, דומה. חלונות, קבצים, אופיס.
  \item \textbf{Motorola Ready For}:Moto דגל תומך, דומה ל-DeX.
  \item \textbf{3. HDR: טכנולוגיית טווח דינמי גבוה}:אופיס קל (דואר, מסמכים), PPT, וידאו. לא ל-Excel כבד או עיצוב.
\end{itemize}

\booksection{ד', חיבור לרכב - הדרך הנכונה}

שימוש בטלפון בנהיגה אסור, אך ניווט, מוזיקה, שיחות נחוצים. חיבור לרכב מאפשר שימוש במסך המכונית:

\begin{itemize}
  \item \textbf{Apple CarPlay}:העדפת iPhone. USB או אלחוטי, מסך המכונית מציג iPhone - ניווט, מוזיקה, טלפון. Siri ללא ידיים.
  \item \textbf{Android Auto}:מקביל ל-iPhone. Google Maps, Spotify. בסין מוגבל.
  \item \textbf{Huawei HiCar}:פתרון Huawei, תמיכה גבוהה. ניווט, מוזיקה, וידאו, מצלמת חנייה.
  \item \textbf{Baidu CarLife}:חוצה פלטפורמות. מעט פונקציות אך תאימות רחבה, לרכבים ישנים.
\end{itemize}

\begin{tcbwarn}{אם הגשת בטעות מידע רגיש, לא מספיק למחוק את הקובץ (הוא עדיין בהיסטוריה), צריך להשתמש ב-}
בכל פתרון, השתמש בשליטה קולית. ניווט ומוזיקה מראש, לא בנהיגה.

\end{tcbwarn}

\booksection{ה', הקרנה בישיבות - בלי בלגן}

ישיבות הן התרחיש השכיח והבעייתי ביותר. ניסיון:

\begin{itemize}
  \item \textbf{הגע מוקדם ובדוק}:אל תחכה לתחילת הישיבה. 10 דקות מראש, בדוק. Wi-Fi נכשל יותר ממה שחושבים.
  \item \textbf{הכן כבל כ-Plan B}:החזק כבל HDMI. כשאלחוטי נכשל, כבל מציל.
  \item \textbf{כבה תצוגת התראות}:בהקרנה, WeChat יופיע על המסך. כבה "תצוגה מקדימה" או הפעל "מצב מצגת".
  \item \textbf{"מצגת" ולא "הפעלת טלפון"}:להצגת PPT, פתח ב-WPS, לא להקרין את כל המסך. מקצועי יותר.
\end{itemize}

\begin{tcbpractice}{פעולה: הקרנה מטלפון לטלוויזיה}
\begin{enumerate}
  \item בטלפון: iPhone ממרכז הבקרה לוחץ על "שיקוף מסך"; אנדרואיד מושך למטה ומחפש "הקרנה אלחוטית".
  \item פתח תמונה או וידאו.
  \item נסה אלחוטי (iPhone: שיקוף, Android: כפתור).
  \item אם נכשל, בדוק תמיכה או נסה DLNA.
  \item אחרי הצלחה, נסה פעולות אחרות, בדוק סנכרון.
\end{enumerate}

\end{tcbpractice}

\begin{tcbsummary}{סיכום פרק}
\begin{itemize}
  \item 3 מערכות: AirPlay, Miracast, DLNA
  \item כבל יציב יותר, חובה בישיבות
  \item DeX/Huawei מצב שולחן = מחשב קל
  \item חיבור רכב: 4 פתרונות
  \item ישיבות: בדוק, הכן Plan B, כבה התראות
\end{itemize}

\end{tcbsummary}
